\documentclass[zh,course]{lecturenotes}
\usepackage{lipsum} % 仅用于示例占位文本
%
% 指定随仓库提供的 Founder 字体，使用相对路径加载
\setnotesCJKmainfont[Path=fonts/Founder/]{FZShuSong-Z01_GB18030 Regular.ttf}
\setnotesCJKsansfont[Path=fonts/Founder/]{FZHei-B01_GB18030 Regular.ttf}
\setnotesCJKmonofont[Path=fonts/Founder/]{FZFangSong-Z02_GB18030 Regular.ttf}
%
% 基本信息
\title{课堂笔记模板示例}
\subtitle{中文模式演示}
\shorttitle{中文示例}
\ccode{代码 7.25}
\subject{示范课程}
\author{作者姓名}
\spemail{speaker@email.com}
\email{email@email.com}
\speaker{授课教师}
\date{12}{06}{2025}
\dateend{19}{06}{2025}
\conference{理学院报告厅}
\place{未来大学}
\flag{示例 \texttt{lecturenotes.cls} v3 中文模式}
\attn{激活 \texttt{zh} 选项后，本类会自动启用 \texttt{ctex}，并可使用 \texttt{\textbackslash setnotesCJKmainfont} 等命令灵活配置字体。}
\morelink{vhbelvadi.com}
%
\begin{document}
\section{课程导入}
本模板的中文模式仅在文档类选项中指定 \texttt{zh} 或 \texttt{chinese} 时才会启用。保持原有英文等模式不变，有助于兼容既有项目。\margintext{边注示例：本段落展示了中文模式下的侧注效果。}

在中文模式下，可以使用标准的数学、定理与表格环境。例如，以下是一个简单的速度公式：
\begin{equation}
\mathbf{v} = \frac{\mathbf{s}}{t} \label{eq:velocity-zh}
\end{equation}%
通过 Lua\LaTeX{} 编译可获得优秀的字形与断行效果。

\lecture[90 分钟]{12}{06}{2025}
在讲次信息中依旧可以记录授课时长、日期等内容。还可以照常插入列表、引用与其他环境。\lipsum[1]

\section{中文字体配置}
\subsection{自定义 CJK 字体}
以下命令在激活中文模式时可用，它们均包装自 \texttt{ctex} 的字体接口：
\begin{description}
  \item[\texttt{\textbackslash setnotesCJKmainfont}] 设置正文中文字体，可接受可选参数传递特定特性，例如字重或 OpenType 特性。
  \item[\texttt{\textbackslash setnotesCJKsansfont}] 设置无衬线中文字体，适合标题或强调文本。
  \item[\texttt{\textbackslash setnotesCJKmonofont}] 设置等宽中文字体，便于展示代码或对齐内容。
\end{description}
这些命令在未启用中文模式时调用会给出提示但不会生效，从而保持旧文档不受影响。

\subsection{更多排版示例}
\lipsum[2]

\begin{table}[h]
\centering
\begin{tabular}{clcc}
\toprule
\multicolumn{2}{c}{试验编号} & 功率 $P_u$ & 正应力 $\sigma_N$\\
\multicolumn{2}{c}{(单位制)} & (W) & (MPa)\\\toprule
\multirow{3}*{甲} & 测试 1 & 285 & 0.26\\\cmidrule(l){2-4}
& 测试 2 & 287 & 0.27\\\cmidrule(l){2-4}
& 测试 3 & 230 & 0.21\\\midrule
\multirow{3}*{乙} & 测试 1 & 430 & 0.19\\\cmidrule(l){2-4}
& 测试 2 & 433 & 0.20\\\cmidrule(l){2-4}
& 测试 3 & 431 & 0.19\\\bottomrule
\end{tabular}
\caption{使用 \texttt{booktabs} 后的示例表格。}
\label{tab:mori-zh}
\end{table}

\lipsum[3]

\section{结构延伸示例}
\subsection{更多讲次}
\lecture[90 分钟]{13}{06}{2025}
\lipsum[4]

\subsection{定理与练习}
\begin{theorem}
若 $f$ 在闭区间 $[a,b]$ 上连续，则其在该区间必定取得最大值与最小值。
\end{theorem}
\begin{exercise}
求解方程 $x^2 - 5x + 6 = 0$ 的全部实根。
\end{exercise}
\begin{remark}
上述结论在更一般的紧致空间中仍然成立。
\end{remark}

\subsubsection{补充内容}
\lipsum[5]

\vskip7ex
\centering
* * *
\end{document}
